% ==============================================================================
% config/pacotes.tex
% Pacotes usados pelo modelo. Em geral, não é preciso alterar este arquivo.
% Se precisar de um pacote novo, acrescente-o ao final, na área indicada.
% ==============================================================================

% ------------------------------------------------------------------------------
% Fontes
% ------------------------------------------------------------------------------
% O modelo compila com pdfLaTeX (padrão), XeLaTeX ou LuaLaTeX. O pacote iftex
% identifica o compilador e carrega apenas o que cada um precisa.
\usepackage{iftex}
\ifPDFTeX
  \usepackage[T1]{fontenc}
  \usepackage{lmodern}   % Latin Modern Mono para códigos (texto e títulos usam Libertine)
\fi

% Linux Libertine (texto) e Linux Biolinum (sem serifa). A opção mono=false
% mantém a fonte monoespaçada padrão, mais legível em listagens de código.
\usepackage[mono=false]{libertine}

% Matemática com o mesmo desenho da Libertine. O libertinust1math já define os
% símbolos do amssymb, por isso este último não deve ser carregado.
\usepackage{amsmath}
\usepackage{libertinust1math}

% Vírgula decimal em modo matemático sem espaço extra: $3,14$.
\usepackage{icomma}

% Ajustes finos de justificação e hifenização.
\usepackage{microtype}

% Recuo também no primeiro parágrafo de cada seção.
\usepackage{indentfirst}

% ------------------------------------------------------------------------------
% Cores, figuras e gráficos
% ------------------------------------------------------------------------------
\usepackage[dvipsnames,table]{xcolor}
\usepackage{graphicx}
\usepackage{subcaption}
\usepackage{pgfplots}
\pgfplotsset{compat=1.18}
\usetikzlibrary{positioning,arrows.meta,shapes.multipart,shapes.geometric,fit,calc,backgrounds}

% ------------------------------------------------------------------------------
% Tabelas
% ------------------------------------------------------------------------------
\usepackage{array}
\usepackage{booktabs}
\usepackage{multirow}
\usepackage{tabularx}
\usepackage{longtable}

% ------------------------------------------------------------------------------
% Códigos e algoritmos
% ------------------------------------------------------------------------------
\usepackage{listings}
\usepackage{algpseudocode}

% ------------------------------------------------------------------------------
% Caixas de orientação e exemplos (podem ser desligadas no arquivo principal)
% ------------------------------------------------------------------------------
\usepackage{comment}
\usepackage[most]{tcolorbox}

% ------------------------------------------------------------------------------
% Citações e referências (ABNT NBR 10520:2023 e NBR 6023)
% ------------------------------------------------------------------------------
% Requer biblatex-abnt 4.0 ou mais recente (TeX Live 2024 em diante) e o Biber
% como processador de bibliografia.
\usepackage[autostyle]{csquotes}
\usepackage[style=abnt, backend=biber]{biblatex}

% ------------------------------------------------------------------------------
% Referências cruzadas e siglas
% ------------------------------------------------------------------------------
\usepackage[nameinlink,noabbrev]{cleveref}
\usepackage{acro}

% ------------------------------------------------------------------------------
% Notas de revisão entre aluno e orientador
% ------------------------------------------------------------------------------
\ifmostrarnotas
  \usepackage[colorinlistoftodos,textsize=footnotesize]{todonotes}
\else
  \usepackage[disable]{todonotes}
\fi

\usepackage{xspace}

% ------------------------------------------------------------------------------
% Pacotes adicionais do aluno
% ------------------------------------------------------------------------------
% \usepackage{...}
